% Espacio fijo por sección para el reporte de la entrega final.
%
% Cada sección va dentro de \begin{espacio}[nombre]{<páginas>} ... \end{espacio},
% con el espacio que le corresponde en páginas (se admiten decimales: 1.55 es una
% página y algo más de media). Al terminar la sección se rellena con espacio en
% blanco lo que le sobra, así que la siguiente empieza siempre en el mismo lugar y
% el PDF tiene siempre las mismas páginas: la suma de los espacios. El nombre, si se
% da, sale en los avisos; si no, sale "sección N".
%
% Si una sección se pasa de su espacio, la compilación lo avisa ("Espacio fijo:
% sección N (... páginas) se pasa por N líneas") y aparece una marca amarilla al
% final de la sección. Ninguna sección se recorta: hay que acortar el texto hasta
% que quepa.
%
% Cómo funciona: al comienzo y al final de cada sección anota en qué página va y
% cuánto lleva ocupado de ella (\pagetotal), y con eso calcula cuánto ocupó la
% sección y cuánto le falta rellenar. Todo en una sola compilación. Antes de medir
% pone un punto de corte neutro (\penalty0) para que TeX decida si la página ya
% se llenó: sin él, \pagetotal puede pasar del alto de la página sin que la página
% haya cambiado todavía.
%
% \espaciovisibletrue escribe en gris, al comienzo del relleno, cuántas líneas le
% sobran a la sección. Para el PDF que se entrega: \espaciovisiblefalse.

\makeatletter
\newif\ifespaciovisible
\espaciovisibletrue
\newdimen\espacio@presupuesto
\newdimen\espacio@usado
\newdimen\espacio@resto
\newdimen\espacio@libre
\newdimen\espacio@inicio
\newcount\espacio@paginaini
\newcount\espacio@excedidas

% Lo que lleva ocupado la página actual, con la profundidad de la última línea
% (0 si está vacía)
\newcommand\espacio@total{\ifdim\pagegoal=\maxdimen 0pt\else\dimexpr\pagetotal+\pagedepth\relax\fi}

\newenvironment{espacio}[2][]{%
  \par\penalty\z@
  \espacio@presupuesto=#2\textheight
  \def\espacio@paginas{#2}%
  \def\espacio@nombre{#1}%
  \espacio@paginaini=\c@page
  \espacio@inicio=\espacio@total
}{%
  \par\penalty\z@
  \ifx\espacio@nombre\@empty\edef\espacio@nombre{sección \thesection}\fi
  \espacio@usado=\dimexpr\numexpr\c@page-\espacio@paginaini\relax\textheight
    - \espacio@inicio + \espacio@total\relax
  \espacio@resto=\dimexpr\espacio@presupuesto-\espacio@usado\relax
  \ifdim\espacio@resto<-0.5pt
    \espacio@excedido
  \else
    \espacio@rellenar
  \fi
}

\newcommand\espacio@rellenar{%
  \edef\espacio@n{\espacio@lineas{\espacio@resto}}%
  \ifdim\pagegoal=\maxdimen
    \espacio@libre=\textheight
  \else
    \espacio@libre=\dimexpr\pagegoal-\espacio@total\relax
  \fi
  \ifespaciovisible
    \ifdim\espacio@resto>\baselineskip
      \nointerlineskip
      \vbox to 0pt{\vskip 0.3\baselineskip
        \centering\footnotesize\color{gris}%
        Espacio libre de esta sección: \espacio@n\ líneas
        de \espacio@paginas\ páginas\par\vss}%
      \nointerlineskip
    \fi
  \fi
  \ifdim\espacio@resto>\espacio@libre
    % No cabe en esta página: se completa la página y se sigue en las siguientes
    \advance\espacio@resto by -\espacio@libre
    \newpage
    \loop\ifdim\espacio@resto>\textheight
      \null\newpage
      \advance\espacio@resto by -\textheight
    \repeat
    % Arriba de la página TeX agrega \topskip antes del relleno: se descuenta
    \advance\espacio@resto by -\topskip
    \ifdim\espacio@resto>0pt \vspace*{\espacio@resto}\fi
  \else
    \ifdim\espacio@resto>0pt \vskip\espacio@resto\relax\fi
  \fi
  \par
}

\newcommand\espacio@excedido{%
  \espacio@resto=-\espacio@resto
  \edef\espacio@n{\espacio@lineas{\espacio@resto}}%
  \global\advance\espacio@excedidas by 1
  \PackageWarningNoLine{espacio}{Espacio fijo: \espacio@nombre\space
    (\espacio@paginas\space páginas) se pasa por \espacio@n\space líneas}%
  \par\noindent
  \colorbox{amarillo}{\parbox{\dimexpr\linewidth-2\fboxsep}{\footnotesize\bfseries
    Esta sección se pasa por \espacio@n\ líneas de su espacio
    (\espacio@paginas\ páginas). Hay que acortarla.}}\par
}

% Una longitud expresada en líneas de texto, redondeada
\newcommand*\espacio@lineas[1]{%
  \the\numexpr\dimexpr#1\relax/\dimexpr\baselineskip\relax\relax}

\AtEndDocument{%
  \ifnum\espacio@excedidas>0
    \PackageWarningNoLine{espacio}{Espacio fijo: \the\espacio@excedidas\space
      secciones se pasan de su espacio}%
  \fi}
\makeatother
